% Reconstructed after the 07:43 UTC execution-environment loss, 16 November 2026.
% The accepted mathematical proof is reconstructed from retained conversation;
% supporting finite checks are rerun against the verified checkpoint-06 inputs.
\subsection*{A dihedral symmetry and the actual cyclic rigid kernel}

\statement{B.65}{Proved}{The pure cyclic centralizer also centralizes $D_5$, and its rigid kernels coincide}
Let $P=P\Gamma_{0,12}$, let $D=\langle d\rangle$ be the chosen
order-five subgroup, and let $H_5=C_P(D)\cong F_6$ as in B.28.
There is an actual involution $r\in T$ such that
\[
 J=\langle d,r\rangle\cong D_5,
 \qquad T\cap J=\langle r\rangle,
 \qquad\langle T,J\rangle=A_5,
                                                               \tag{B.65.1}
\]
and
\[
 C_P(D)=C_P(J)=H_5,\qquad
 C_{\widehat P}(D)=C_{\widehat P}(J)=\widehat H_5.
                                                               \tag{B.65.2}
\]
The two rigid-seven forgetful maps whose sets are $D$-stable
vanish on $H_5$.  The other ten maps have exactly the same
kernel on $H_5$, and exactly the same kernel on
$\widehat H_5$, after their specified target transports.
This common kernel is the kernel of an explicit six-generator
homomorphism
\[
 \theta_5:F(f_0,\ldots,f_5)\longrightarrow PB_6,
                                                               \tag{B.65.3}
\]
and its completed analogue is $\ker\widehat\theta_5$.
No injectivity of (B.65.3), nor identification of the latter
kernel with the embedded closure of the former, is asserted.

\begin{proof}
Use the actual cyclic marking of B.60--B.61: labels $1,3$
are $0,\infty$, and the two five-cycles are
\[
 X=(2,6,12,9,4),\qquad Y=(5,7,11,10,8).
\]
Put $X_i=\zeta^i$ and $Y_i=u\zeta^i$, where
$u\in S_5=\mathbb P^1\setminus(\{0,\infty\}\cup\mu_5)$.
At the common icosahedral point, $u=-\varphi^2$.
The involution
\[
                    r_u(z)=u/z
                                                               \tag{B.65.4}
\]
preserves the twelve-point configuration, exchanges its two
poles, and sends $X_i$ to $Y_{-i}$.
Its permutation in the common marking is
\[
       r=(1\ 3)(2\ 5)(6\ 8)(12\ 10)(9\ 11)(4\ 7).
                                                               \tag{B.65.5}
\]
It satisfies $rdr=d^{-1}$.  The exact permutation check
against the twelve tetrahedral permutations proves
$r\in T$, $|J|=10$, $T\cap J=\{1,r\}$ and
$\langle T,J\rangle=A_5$ of order sixty.  In particular
the involution is the actual common group element, not
an abstract additional symmetry.

Here is the discrete commuting argument.
Let $u(t)$ be a based loop in $S_5$ and choose an ambient
isotopy $g_t$ carrying the labelled base configuration to
the configuration at $u(t)$, with $g_0=1$.
The path
\[
                    g_t^{-1}r_{u(t)}g_t
\]
preserves the base configuration and has the fixed
permutation (B.65.5).  It joins $r$ to $h^{-1}rh$,
where $h$ is the monodromy mapping class.
Consequently $h$ commutes with $r$ in the full mapping-class
group.  Every element of $H_5$ is such a monodromy by
B.28, so $H_5\subset C_P(J)$, and the reverse inclusion
follows from $D\leq J$.
For the completed statement, B.28 identifies
$C_{\widehat P}(D)$ with the actual embedded closure
$\widehat H_5$.  Commutation with the finite element $r$
is a closed condition, so the entire closure still commutes
with $J$.  This proves (B.65.2) without interchanging an
unknown intersection with completion.  Thus one may also
write the actual residual centralizer as
\[
       \mathcal K=C_{\widehat P}(T)\cap C_{\widehat P}(J).
                                                               \tag{B.65.6}
\]

The two $D$-stable seven-sets are
\[
 S_X=\{1,2,3,4,6,9,12\},\qquad
 S_Y=\{1,3,5,7,8,10,11\}.
\]
Their labelled moduli maps are constant: the first family
is fixed, and the second is made fixed by $z\mapsto z/u$.
Hence both forgetful homomorphisms on $H_5$ are zero.
Set
\[
 S_+=\{1,2,4,5,6,10,11\},\qquad
 S_-=\{2,3,4,7,10,11,12\}.
\]
The remaining ten sets form one regular $J$-orbit.
The precise transports are as follows; the last column
records the image of the label $2$.
\[
\begin{array}{c|c|c}
 a&aS_+&a(2)\\\hline
 1&\{1,2,4,5,6,10,11\}&2\\
 d&\{1,2,6,7,8,10,12\}&6\\
 d^2&\{1,5,6,8,9,11,12\}&12\\
 d^3&\{1,4,5,7,9,10,12\}&9\\
 d^4&\{1,2,4,7,8,9,11\}&4\\
 r&\{2,3,5,7,8,9,12\}&5\\
 dr&\{3,4,5,6,7,9,11\}&7\\
 d^2r&\{2,3,4,7,10,11,12\}&11\\
 d^3r&\{2,3,6,8,9,10,11\}&10\\
 d^4r&\{3,4,5,6,8,10,12\}&8
\end{array}
                                                               \tag{B.65.7}
\]
For an actual $h\in H_5$, equivariance of forgetting and
the already proved equality $aha^{-1}=h$ give
\[
          f_{aS_+}(h)=a_*f_{S_+}(h).
                                                               \tag{B.65.8}
\]
Here $a_*$ is the isomorphism of target mapping-class
groups induced by that labelled sphere homeomorphism.
The same identity holds on the completions by continuity.
This proves equality of the ten kernels in both categories.
The rigid five- and six-point forgettings factor through
appropriate rigid seven-point forgettings, as in B.13.

For a completely specified finite-word test, normalize
labels $2,4,6$ to $\infty,0,1$ in the actual common
marking of B.61.  The finite labels of $S_+$ occur in order
\[
                       (4,11,5,1,10,6).
                                                               \tag{B.65.9}
\]
Delete the other five finite strands from each of the six
original cyclic eleven-strand words in
\nolinkurl{next_B59_image/common_base_braid_words.json}.
The resulting words define (B.65.3); their unreduced lengths
are $36,34,22,30,54,32$, respectively, and their exponent
sums are $12,-10,4,-8,-8,4$.
The reconstruction check under
\nolinkurl{next_joint_rigid_symplectic/}
stores all six words, verifies their purity and their
linking numbers, and independently compares the deletion
with the restored checkpoint-six rigid-restriction data.

Let $z=\Delta_6^2$ be the central full twist.
Capping gives $PB_6/\langle z\rangle=P\Gamma_{0,S_+}$.
The two retained points $4,6$ are fixed throughout, so
the pair-linking homomorphism $\ell_{4,6}:PB_6\to\mathbb Z$
vanishes on all six source images, whereas
$\ell_{4,6}(z)=1$.
It follows that the capped kernel is exactly
$\ker\theta_5$.
This remains true with the full profinite topology:
the central splitting by $z$ gives
$PB_6=\ker(\ell_{4,6})\times\langle z\rangle$,
and hence the corresponding direct product after completion.
The intersection of $\ker\widehat\ell_{4,6}$ with the
completed central factor is trivial.
Thus the capped completed kernel is
$\ker\widehat\theta_5$ as claimed.
For a finite source word, its membership in
$\ker\theta_5$ is decided by substitution of these six
braids followed by their faithful Artin action on a free
group of rank six.  This is a kernel-membership test;
it is not a general subgroup-membership algorithm.

Finally, equality of the full target kernels does not
identify every further quotient of the two targets.
The transport $d^2r:S_+\to S_-$ sends $2$ to $11$.
The genus-two double-cover homology test that omits
label $2$ on each of $S_+,S_-$ therefore omits different
points after transport.  The two selected finite homology
tests can consequently be independent even though the
full seven-point kernels agree.  No conclusion about
injectivity of the remaining full map follows from either
finite homology calculation.
\end{proof}
